\documentclass[11pt,a4paper]{ltjsarticle}
\usepackage[margin=22mm]{geometry}
\usepackage{booktabs,tabularx,array,enumitem,xcolor}
\usepackage{tikz}
\usetikzlibrary{positioning,arrows.meta}
\usepackage[hidelinks]{hyperref}
\setlist{itemsep=1pt,topsep=3pt,parsep=0pt}
\newcolumntype{L}{>{\raggedright\arraybackslash}X}
\newcolumntype{P}[1]{>{\raggedright\arraybackslash}p{#1}}
\renewcommand{\arraystretch}{1.25}
\newcommand{\voice}[1]{\begin{quote}\small 「#1」\end{quote}}

\title{MER-DA の組み直しの計画\\[2pt]\large 何をしたいか，何ができているか，何ができていないか，何をするか}
\author{2026-11-04 の相談のまとめ}
\date{2026年11月4日}

\begin{document}
\maketitle

\section{この文書について}

2026-11-04 の研究者との相談の内容を，1 つの文書にまとめた．
今後の設計と実験は \texttt{MER-DA/DAI/} で行う．
数値の後に付けた「照合済み」は，検査役が結果ファイルと値を照らしたものである．
「検査前」は，2026-10-28 に検査の仕組みを入れる前に記録した値で，論文に書く前に照合が要る．
用語は最後の節にまとめた．

\textbf{第 2 版（2026-11-04）}：ML の査読者を想定した役が第 1 版を検査し，「条件つき OK」と判定して重大な指摘を 8 つ挙げた（\texttt{DAI/checks/計画\_2026-11-04/review\_ml.md}）．その対応を 8 節「設計の約束」に書き，7 節の順番を直した．

\section{何をしたいか}

研究者の言葉（2026-11-04）を，そのまま引く．

\voice{このAlbanieの研究を大いに参考にしてMERーDAを組み直したい．現在Albanieよりも以下の点で複雑であるが拡張版をやりたい．複雑な版．（k個のラベル付きデータ，モダリティが３つ．よくさせたいのは各単一モダリティと３つの融合，ソースとターゲットで別コーパスのDA設定，などなど）}

\voice{MERー DAでやりたいことは，Albineないしその後続研究を見て，それらの発展．特にMERにおいてマルチモダリティせいをどう使うのがいいかを調べて，アルゴリズムやモデルを作る}

\voice{DAIをターゲットに，採択されることを目指したい}

\voice{俺は知りたいんだ，何が起きているかを．つまり説明をしたい．またそこから提案手法を提案したい}

まとめると，目的は次の 3 つである．
\begin{enumerate}
  \item \textbf{説明する．}別のモダリティの予測をラベル無しの発話に付けて学ぶと，ターゲットの識別器がなぜ，どこで良くなるのかを説明する．
  \item \textbf{手法を作る．}その説明から，識別器の作り方（アルゴリズムやモデル）を導く．
  \item \textbf{DAI に採択される．}投稿先は Discover Artificial Intelligence（Springer Nature）である．DAI の査読は研究の重要さを評価せず，結果が正確に報告されているかと，研究に根本的な欠陥が無いかを見る（投稿規定の Step 3．\texttt{DAI\_paper/docs/venue.md}）．
\end{enumerate}

\section{設定}

Albanie ほか（2018）は，顔の表情の識別器の予測を，感情のラベルが無い VoxCeleb の音声に正解として付け，音声の識別器を学んだ．
この研究では，それを次の形に広げる．

\begin{center}
\begin{tabularx}{\linewidth}{@{}p{30mm}L@{}}
\toprule
項目 & 中身 \\
\midrule
コーパス & AFEW，IEMOCAP，MOSEI，MELD．音声・顔・言葉の 3 つがそろう 4 コーパスである．AFEW の言葉は，2026-11-04 に Whisper large-v3 で書き起こして足した． \\
分け方 & 1 つのコーパスをターゲットにし，残りの 3 つをソースにする．4 つを順にターゲットにする（leave-one-corpus-out）． \\
ソースの資産 & ソースのコーパスごとに，音声・顔・言葉の識別器と，3 つをつないだ融合（AVT）の識別器を作り，保存して後で使う． \\
ターゲットの資産 & 感情ごとに K 本のラベル付きの発話と，ラベルの無い発話（3 つのモダリティがそろう）． \\
良くしたい識別器 & ターゲットの音声・顔・言葉の各識別器と，AVT の融合の識別器．どれも，凍結した事前学習の特徴の上の線形の識別器（ロジスティック回帰）である． \\
評価 & UAR（感情ごとの正解率の平均．でたらめなら 25\%）．ターゲットの中では，テストの話者を学習側に入れない． \\
比べる条件 & メタ学習あり・なし．比較手法． \\
\bottomrule
\end{tabularx}
\end{center}

\begin{center}
\begin{tikzpicture}[
  box/.style={draw,rounded corners,align=center,font=\small,minimum height=10mm,text width=40mm},
  arr/.style={-{Stealth[length=2.2mm]},thick}]
\node[box] (src) {ソースの 3 コーパス\\（ラベル付き）};
\node[box,right=12mm of src] (clf) {コーパスごとの識別器\\音声・顔・言葉・AVT};
\node[box,below=9mm of clf,text width=52mm] (pl) {ターゲットのラベル無しの発話に\\予測を付ける（疑似ラベル）};
\node[box,left=12mm of pl] (k) {ターゲットの K 本\\（ラベル付き）};
\node[box,below=9mm of pl] (stu) {ターゲットの識別器を学ぶ\\音声・顔・言葉・AVT};
\draw[arr] (src) -- node[above,font=\scriptsize]{学習} (clf);
\draw[arr] (clf) -- (pl);
\draw[arr] (pl) -- (stu);
\draw[arr] (k) |- (stu);
\end{tikzpicture}
\end{center}

\section{先行研究で分かっていること}

2026-11-04 に，Albanie らを引用した論文（Semantic Scholar 320 件，OpenAlex 314 件．重なりを除いて約 325 件）と，関連分野を調べた．
報告は \texttt{research\_notes/先行研究\_Albanieの後続研究\_2026-11-04/} にある．

\subsection{結論}

拡張の 4 点（ターゲットに K 本のラベル，3 モダリティ，各単一モダリティと融合の両方を良くする，ソースとターゲットが別コーパス）を同時に扱った研究は，見つからなかった．
4 点のそれぞれには先例がある．

\subsection{既に報告されていること}

下の 1〜3 の数値は，原文の文字で確かめた．4 は調査役が原文で確かめた．
\begin{enumerate}
  \item \textbf{Albanie らの方法をそのまま別コーパスに当てると弱い．}Zhang ほか（2021）では，EmoVoxCeleb で学んだ音声の識別器を IEMOCAP の 4 感情に当てると UA 27.99\% だった．IEMOCAP の全ラベルで微調整しても 54.97\% で，ラベルだけで学んだ 59.43\% を下回った．顔と音声の予測を混ぜた疑似ラベルでは 62.50\% だった．
  \item \textbf{別のモダリティの疑似ラベルの得は，ラベルが増えると減り，負にもなる．}ACL-MER（2025）の MER2023 では，正解率が，ラベル 100 本のとき 0.432 から 0.674 に上がり，1000 本のとき 0.826 から 0.819 に下がった．
  \item \textbf{2 つのモダリティで予測を付けると，1 つより正確である．}SetPeER（2025）では，CMU-MOSEI の人のラベルとの一致率が，音声 52.2\% から音声＋言葉 57.2\% に，顔 37.4\% から顔＋言葉 50.8\% に上がった．
  \item \textbf{弱いモダリティの識別器の予測は，学ぶ側を悪くする．}TelME（2024）の IEMOCAP では，言葉の識別器が，音声の識別器の予測を学ぶと 66.57 から 57.42 に下がった．MST-Distill（2025）の RAVDESS では，顔の識別器が，音声の識別器の予測を学ぶと下がり，融合の識別器の予測を学ぶと上がった．
\end{enumerate}

\subsection{見つからなかったこと}

\begin{itemize}
  \item 別のモダリティの疑似ラベルの得を，K を 0 から数十本まで変えて測った研究．
  \item ソースだけで学んだ識別器を，ターゲットのラベル無しの発話に当てた研究．後続研究は，第三のコーパスか，同じ収集元のデータに当てている．
  \item 3 つの単一モダリティと融合を同じ方法で学び直し，全部の成績を出した研究．
  \item 同じ識別器に，予測を付ける側の種類（別のモダリティ 1 つ，別のモダリティ 2 つの融合，学ぶ側のモダリティを含む融合）を並べて比べた研究．
  \item 疑似ラベルがどの感情の対の誤りを直すかを測った研究．
\end{itemize}

研究者の指摘のとおり，見つからなかったことの多くは現象か実験の設定であり，「なぜそうなるか」の説明を含まない．このままでは論文の主張として弱い．

\section{MER-DA でできていること}

\subsection{分かっていること}

出典は \texttt{DAI\_paper/inventory.md} である．「位置の誤り」と「方向の誤り」は最後の節で定義した．\textbf{この表の値は前の設定（7 群，全ソースをまとめた出発点，対話を除かない分割）で測ったもので，研究の方向を考える材料にだけ使う．論文に書く値は，すべて新しい設定の出力から計算し直す}（8.6 節）．

\begin{center}\small
\begin{tabularx}{\linewidth}{@{}LP{40mm}P{20mm}@{}}
\toprule
分かっていること & 数値 & 状態 \\
\midrule
K 本は位置の誤りを安く直すが，方向の誤りは遅くしか直さない & 音声では K=10 でも方向の誤りの 6〜7 割が残る & 照合済み（一部は検査前） \\
ソースのデータを K 本と混ぜて学び直すと，K が大きいとき損になる & 音声の K=25 で $-1.69$ & 照合済み \\
同じモダリティのラベル無しの発話は，位置を直すが，方向を直さない & ラベル無しの 8 手法の 210 単位のうち，方向の誤りを減らしたのは 7（偶然なら約 5），増やしたのは 79 & 照合済み \\
別のモダリティを含む識別器の予測は，方向を直す & 81 単位のうち，減らしたのは 38，増やしたのは 1 & 照合済み \\
予測を付けるのは，学ぶ側のモダリティを含む融合がよい & 音声の識別器に対し，融合は $+6.7$〜$+8.7$，顔だけは $-4.2$ & 検査前 \\
予測の形（柔らかい，温度つき，硬い，釣り合いを取る）を替えても差が出ない & 21 の組のうち 0 & 照合済み \\
学ぶ側は，予測を付ける識別器が勝つ分の約 0.45 しか受け取らない & 勝つ幅 $+7.3$〜$+8.0$ に対し得 $+3.4$ & 検査前 \\
発話ごとに正しいモダリティを選べれば正解率は高いが，ラベル無しでは選べない & 選べれば 74\%，融合 61.2\%，ラベル無しの選び方 52.0\% & 検査前 \\
ソースごとの識別器の予測をまとめると，全ソースをまとめた 1 つの識別器より当たる & まとめ方を仮のターゲットで学ぶと，音声 $+3.3$，顔 $+4.6$，音声＋顔 $+4.8$ & 検査前 \\
設定を仮のターゲットでまとめて選ぶと効く & K=5 で $+1.0$〜$+1.6$ & 検査前 \\
\bottomrule
\end{tabularx}
\end{center}

\subsection{使える道具}

\begin{itemize}
  \item 7 群の凍結した特徴．各モダリティにエンコーダが 2 つずつある（音声 WavLM・HuBERT，顔 ResNet-50・CLIP，言葉 RoBERTa・DeBERTa）．
  \item 話者を分ける分割，K 本を入れ子に抜く手順，GPU で解くロジスティック回帰．
  \item 「K 本＋疑似ラベル＋出発点への L2 正則化」で学び直す学習器．
  \item 比較手法の実装．ラベル無しの標準手法（自己学習，エントロピー最小化，IM，T3A，ラベル伝播），温度つき蒸留，共訓練，共正則化，DANN，MME，xMUDA 型，CORAL．どれも全ソースをまとめた出発点を前提にしているので，コーパスごとの識別器に合わせて手直しが要る．
  \item ソースごとの識別器の予測をまとめる 6 通りの方法（多数決，まとめ方を仮のターゲットで学ぶ方法，Dawid--Skene など）．
  \item 位置と方向の誤りを分ける計算．実験の検査役（関門）の仕組み．RunPod の H100 を使う設定．
\end{itemize}

\section{できていないこと}

\subsection{説明が無いもの}

\begin{itemize}
  \item \textbf{積のモデルが強い理由．}積のモデルは，ターゲットの予測で各モダリティの識別器を学び直し，3 つの確率を掛けたものである．仮のターゲットで設定を選ぶと，81 回のうち 77 回でこれが選ばれた（照合済み）．ソースだけの融合を上回る理由は分かっていない．
  \item \textbf{学ぶ側が約 0.45 しか受け取らない理由．}
  \item \textbf{顔の識別器の得が小さい理由．}顔は出発点が最も弱いのに，得が最も小さい．
  \item \textbf{予測の価値が K とともにどう変わるか．}検査を通した記録が無い．
  \item \textbf{どの予測がどの感情の対の方向を直すか．}この解析は設計済みで，検査役の直しを入れてから計算する．
\end{itemize}

\subsection{設定の不備}

\begin{itemize}
  \item \textbf{資産の定義．}適応の段でソースのデータを使ってよいかを決めていない（7.2 節）．
  \item \textbf{IEMOCAP の同じ台詞．}IEMOCAP の台本では 10 人全員が同じ台詞を読む．即興と台本を 1 つにしてセッションを 1 つテストにすると，学習側の 10.4〜12.0\% がテストと全く同じ書き起こしを持つ（台本では 14〜16\%，即興では 7〜8\%）．
  \item \textbf{テストと候補が同じ対話を共有する．}候補は，ターゲットのラベル無しの発話である．前回の分割では，IEMOCAP 即興のテストの 99.5\%，MELD の 90.3\% が候補と同じ対話に属した．
  \item \textbf{評価に使うコーパスで手法の形も決めた．}仮のターゲットだけで決め直す計算は済んだが，照合はまだである．
  \item \textbf{ターゲットが 4 つしかない．}統計の単位が少ない．
  \item \textbf{同じモダリティの手法との比較は，使う情報の比較である．}同じモダリティの手法は，別のモダリティの予測を受け取らない．
\end{itemize}

\subsection{検査が終わっていないもの}

\begin{itemize}
  \item 仮のターゲットだけで手法の形を決め直した計算：済み，照合はまだ．
  \item テストと同じ対話を候補から除いた比べ直し：照合の 1 項目の許容幅が研究者の返事待ち．
  \item 感情の対ごとの分解：設計済み，未計算．
\end{itemize}

\section{何をするか}

\subsection{決めたこと（2026-11-04）}

\begin{enumerate}
  \item 作業は \texttt{MER-DA/DAI/} で行う．このフォルダの規則は \texttt{DAI/CLAUDE.md} に書く．
  \item AFEW，IEMOCAP，MOSEI，MELD の 4 コーパスで leave-one-corpus-out の実験をする．
  \item ソースのコーパスごとに，音声・顔・言葉と AVT の識別器を作って保存する．
  \item メタ学習あり・なしの両方を回す．
  \item 勾配法で学ぶ比較手法（DANN，MME，xMUDA 型，CORAL）も回す．前回の同じ規模の計算は，H100 の計算時間の合計で 8.4 時間だった．RunPod を使う．
  \item 比較手法は，最も近い先行研究が何と比べていたかを見て決める．
  \item 実施する実験は，計算の前に ML の査読者と感情認識の査読者を想定した役の両方から OK を得る．査読想定の役は，ML の査読者，感情認識の査読者，編集者の 3 役にした（\texttt{DAI/.claude/agents/}）．
\end{enumerate}

\subsection{これから決めること}

\begin{center}\small
\begin{tabularx}{\linewidth}{@{}P{28mm}P{42mm}L@{}}
\toprule
決め事 & 推す既定 & 理由 \\
\midrule
IEMOCAP の扱い & 即興だけを使う & 同じ台詞の問題が，特別な除外なしに消える．書き起こしが全く同じ発話を除いても，少しだけ違う台詞は残る．台本の 2,588 本を失う．ソースのときも即興だけを使い，コーパスの中身をソースとターゲットでそろえる（8.8 節）． \\
資産の定義 & 適応の段ではソースのデータを使わず，識別器だけを使う．識別器を作る段とメタ学習の練習では，ソースのデータを使う & ソースのモデルだけを使う適応は，source-free DA（Liang ほか 2020 の SHOT）や仮説転移学習（Kuzborskij \& Orabona 2013）と同じ形で，先行研究にある設定である．DANN，MME，CORAL は，適応の段でもソースのデータを使う相手として並べ，使う情報の違いを表に書く． \\
識別器の罰則の強さ C & 各コーパスの中で，話者を分けた交差検証で選ぶ & ターゲットのラベルを使わない． \\
平均を引くか & 各コーパスの平均を引いてから学ぶ & ターゲットの平均を引くだけで，176 組の 95\% が損をしなかった（検査前）． \\
対話の除外 & テストと同じ対話の発話を，K 本と候補から最初から除く & 前回の模擬査読の指摘を最初から避ける． \\
融合の種類 & AVT のほかに AV，AT，VT も作る & 計算はどれも数秒である． \\
\bottomrule
\end{tabularx}
\end{center}

\subsection{順番と所要の見込み}

ML の査読者の指摘 M1 を受けて，手法を作る段と評価する段を分けた．評価のターゲットのテストは，手法の形を固定した後の 1 回だけ使う．所要は前回の比較実験の記録から割り戻した見込みで，測ってはいない．

\begin{center}\small
\begin{tabularx}{\linewidth}{@{}P{6mm}LP{42mm}@{}}
\toprule
順 & すること & 見込み \\
\midrule
1 & 決め事を決め，設計を書き，台本を関門の検査役と ML・感情認識の査読者に通す & 作業で数時間 \\
2 & コーパスごとの識別器を作って保存する & GPU で 1 時間以内 \\
3 & 仮のターゲットの上で，全手法を全部の格子で回す & H100 で数時間 \\
4 & 仮のターゲットの出力だけで，設計のための解析をする．評価のターゲットのテストは見ない & 数分〜数時間 \\
5 & \textbf{手法の形，格子，選び方の規則を文書に書いてコミットし，固定する．}\texttt{docs/seen-data.md} に 1 行足す & 作業で 1 時間 \\
6 & 評価のターゲットで全手法を 1 回だけ回す．第 1 波（疑似ラベルの手法とラベル無しの手法）と第 2 波（DANN，MME，xMUDA 型，CORAL） & H100 で約 6 時間＋約 8〜9 時間（前回は 6.1 時間と 8.4 時間） \\
7 & 評価のターゲットの出力で，説明の解析をする．固定の後に形を変えた版は探索として別に書く & 保存した出力から \\
8 & 結果を書く前に照合の検査役に通す & 各回 1 時間ほど \\
\bottomrule
\end{tabularx}
\end{center}

\section{設計の約束（ML の査読者の指摘への対応）}

指摘の番号は \texttt{review\_ml.md} の番号である．ここに書いた既定のうち，研究者がまだ確かめていないものは 8.9 節にまとめた．

\subsection{手法を作るデータと評価するデータ（M1）}

\begin{itemize}
  \item 4 コーパスは leave-one-corpus-out でどれもほかの回の仮のターゲットになるので，設計に使っていないデータは無い．
  \item 7.3 節の順 5 で，手法の形，格子，選び方の規則を固定してコミットする．評価のターゲットでは，その後に 1 回だけ回す．固定の前の設計の解析は，仮のターゲットの出力だけで行う．
  \item 別のコーパスは用意しない．要旨と限界の節に「手法は評価に使ったのと同じ 4 コーパスの上で作った」と書き，比較の結果を確認の結果としては書かない（2026-11-04 研究者の承認）．
\end{itemize}

\subsection{統計（M2）}

\begin{itemize}
  \item \textbf{分割と反復}：各コーパスを 8.5 節の塊で 5 つに分け，5 つを順にテストにする．K 本の抜き方の種を 2 つにし，1 つのターゲットで 10 回にする．IEMOCAP 即興ではテストの集合は 5 通りである．
  \item \textbf{K}：0，1，5，10，25 を計算する．主な比較に使うのは 0，5，25 である．IEMOCAP 即興の学習側の感情ごとの最少は 199 本，AFEW は約 140 本で，K=25 は取れる．
  \item \textbf{エンコーダ}：主は音声 WavLM，顔 ResNet-50，言葉 RoBERTa である．主な比較だけを，HuBERT，CLIP，DeBERTa で繰り返す．
  \item \textbf{区間}：コーパスの中の区間は，テストの塊を引き直すブートストラップで作る．塊の実効の数が 20 未満のときは，塊を 1 つずつ抜く jackknife の区間と比べて広いほうを使う．
  \item \textbf{コーパスをまたぐ文}：「4 コーパスのうち x で」と数で書き，p 値を付けない．
  \item \textbf{主な比較}：コーパスごとに次の族を前もって決め，族の中で Holm の補正をかける．族の外は記述とする．
  \begin{itemize}
    \item 族 1：生徒 4 種（音声，顔，言葉，AVT）× K 3 つで，提案（同時に解く版）と比較の基準（K 本だけ）の差．12 個．
    \item 族 2：生徒 4 種 × K 2 つ（5，25）で，役割を分けた版と同時に解く版の差．8 個．
    \item 族 3：生徒 4 種 × K 3 つで，提案と，同じ情報を受け取る比較相手のうち仮のターゲットで最も良かったものとの差．12 個．
    \item 族 4：生徒 4 種 × K 3 つで，選択あり（仮のターゲットで選ぶ）と選択なし（固定の規則）の差．12 個．
  \end{itemize}
  \item \textbf{同じ水準}：差が幅 $\delta$ の中に区間ごと収まるときだけ「同じ水準」と書く．$\delta$ は，比較の基準の UAR の反復間の揺れを仮のターゲットで測り，その 2 倍の標準誤差とする（前回と同じ規則）．
  \item \textbf{検出できる差}：上の規則で，コーパスごとに主な比較で検出できる差の大きさを，評価のターゲットを回す前に仮のターゲットから計算して書く．
\end{itemize}

\subsection{説明に使う物差し（M3）}

\begin{itemize}
  \item \textbf{位置と方向の誤り}：模型に頼らない次の式を主にする．UAR$(W,b)$ は重み $W$，切片 $b$ の識別器の UAR である．$b^{*}(W)$ は，$W$ を固定して評価の発話のラベルで切片だけを当て直した値，$W^{*}$ は候補の全ラベルで学んだ識別器の重みである．
  \begin{center}
  位置の誤り $=$ UAR$(W, b^{*}(W)) -$ UAR$(W,b)$，\quad 方向の誤り $=$ UAR$(W^{*}, b^{*}(W^{*})) -$ UAR$(W, b^{*}(W))$
  \end{center}
  2 つの和は UAR$(W^{*}, b^{*}(W^{*})) -$ UAR$(W,b)$ に一致する．$b^{*}$ を UAR の最大化で求めるか交差エントロピーで求めるかは，設計の文書で決める．テストのラベルを使う量なので，記述にだけ使い，手法の設定の選択には使わない．感情の対ごとの版も同じ形で定める．前回の 2 つの正規分布の模型の量は，補助として並べる．
  \item \textbf{順序}：方向を先に入れ替える順でも測り，2 つの順序の差を報告する．
  \item \textbf{妥当性の確かめ}：評価のターゲットを使う前に，仮のターゲットで，切片だけを人工的にずらした識別器と，重みだけを回した識別器を作る．前者が位置の誤りだけとして，後者が方向の誤りだけとして測られるかを確かめる．
  \item \textbf{ラベル何本ぶん}：主張には使わない．記述として出すときは，学習曲線を候補の全ラベルの点まで延ばし，読みが曲線の外に出たら「読めない」と書く．
  \item \textbf{代わりの説明}：説明ごとに，立てる説明，代わりの説明，2 つを分ける比較を，計算の前に書く．
  \begin{center}\small
  \begin{tabularx}{\linewidth}{@{}P{30mm}LL@{}}
  \toprule
  立てる説明 & 代わりの説明 & 分ける比較 \\
  \midrule
  K 本は位置，別のモダリティの予測は方向を直す & K 本も方向を主に直している & 役割を分けた版と同時に解く版．K 本で切片だけを直す版と K 本で全部を学び直す版 \\
  積のモデルが強いのは，ターゲットの疑似ラベルで学び直すから & 複数の識別器の確率を掛けること自体が効く & 学び直さない積，学び直した各識別器 1 つずつ，学び直した積を同じ表に並べる \\
  生徒が受け取るのは，予測のうち自分の特徴で再現できる部分だけ & 予測を付ける側の当たり方だけで得が決まる & 予測を付ける側の種類ごとに，当たり方と再現できる度合いの両方で得を説明し，どちらが得をよく説明するかを比べる \\
  顔の得が小さいのは，顔の特徴で再現できる部分が少ないから & 顔の出発点の方向がすでに最良に近く，直す余地が少ない & 顔の生徒の，直す余地（全ラベルの識別器との差）と再現できる度合いを並べる \\
  \bottomrule
  \end{tabularx}
  \end{center}
\end{itemize}

\subsection{比較の公平さと比較相手（M4，M5）}

\begin{itemize}
  \item \textbf{出発点}：全手法に同じ出発点を渡す．出発点は，ソースのコーパスごとの識別器の点数（ソフトマックスの前の値）を等しい重みで平均したものである．これは線形の識別器になる．ソースの名前ごとの重みを持たないので，ソースが 2 つの仮のターゲットの回にも同じ形で作れる（軽微な指摘 m2）．
  \item \textbf{選び方}：全手法の設定を，同じ仮のターゲットの手順で選ぶ．候補の数は手法の間でおよそ 12 通りにそろえる．選択なしの版は，全手法に共通の規則「各つまみの格子の中央（偶数個なら小さいほうの中央）」で決める．提案では，疑似ラベルの重みの合計 500，K 本の倍率 4 になる．
  \item \textbf{出発点への引き戻しの強さ}：比較の基準（K 本だけ）について (生徒，K) ごとに仮のターゲットで選び，全手法で同じ値を使う（軽微な指摘 m3）．
  \item \textbf{取りこぼし}：選んだ設定とテストで最良の設定の差を，手法ごとに出す．
  \item \textbf{書き方}：同じ情報を受け取らない相手との差は「情報の差」として書く．手法の優劣は，同じ情報を渡した相手との差だけで書く．
  \item \textbf{資産の書き方}：「ソースの側で，コーパスごとの識別器と選んだ設定を作って渡す．適応の段ではソースのデータを使わない」と書き，source-free とだけは書かない．
\end{itemize}

\begin{center}\scriptsize
\begin{tabularx}{\linewidth}{@{}P{34mm}P{13mm}P{13mm}P{14mm}P{9mm}P{9mm}P{13mm}L@{}}
\toprule
手法 & 適応で使うソースのデータ & 設定の選択に使うデータ & 別のモダリティの予測 & K 本 & ラベル無し & 入れるか & 理由 \\
\midrule
出発点のまま，K 本だけ（比較の基準），K 本だけでゼロから & なし & 仮のターゲット & なし & 使う & なし & 入れる & 下限の参照 \\
ソースのデータと K 本を混ぜて学ぶ & あり & 仮のターゲット & なし & 使う & なし & 入れる & 半教師ありの DA の標準の基準 \\
候補の全ラベルで学ぶ & なし & なし & なし & 全部 & 全部 & 入れる & 上限の参照 \\
自己学習，エントロピー最小化，IM，T3A，ラベル伝播，割合の補正 & なし & 仮のターゲット & なし & 使う & 使う & 入れる & 同じモダリティのラベル無しの手法 \\
提案（同時に解く，役割を分ける），温度つき蒸留，一致した発話だけの疑似ラベル，別のモダリティと自分の予測を混ぜる手，共訓練，共正則化，xMUDA 型 & なし & 仮のターゲット & あり & 使う & 使う & 入れる & 同じ情報を受け取る相手 \\
DANN，MME，CORAL & あり & 仮のターゲット & なし & 使う & 使う & 入れる & 適応でもソースのデータを使う相手 \\
SHOT & なし & — & なし & — & 使う & 入れない & 特徴の抽出器を学ぶ手法で，特徴を凍結するこの設定では前提が満たされない \\
FixMatch，$\pi$ モデル，Temporal ensembling & なし & — & なし & — & 使う & 入れない & 入力の変形を前提にし，凍結した平均の特徴には変形が無い \\
Co-teaching & なし & — & なし & — & 使う & 入れない & 前回，線形の頭では勾配法が比較の基準の解に届かなかった \\
\bottomrule
\end{tabularx}
\end{center}

予測を付ける側の種類（別のモダリティ 1 つ，別の 2 つの融合，生徒のモダリティを含む融合，学び直さない積，学び直した積，コーパスごとの識別器）は，提案の中の比べとして全部を並べる．

\subsection{テストと学習側の重なり（M6，M7）}

\begin{itemize}
  \item \textbf{塊}：IEMOCAP 即興はセッション（5 個），AFEW は俳優と作品のつながりでまとめた塊（52 個．最大 156 発話＝19.6\%），MOSEI はチャンネル（分からなければ動画），MELD は話者の群とする．
  \item \textbf{分け方}：IEMOCAP，AFEW，MOSEI は塊ごとに 5 つに分ける．MELD は主役が多くの対話に出るので塊では分けず，話者で分けたうえで，テストの対話の発話を K 本と候補から除く．
  \item \textbf{除いた後の数}：除いた後の候補の数と感情ごとの最少の本数を，コーパスごとに数えてから K の上限を確かめる．
  \item \textbf{感度}：MOSEI は，チャンネルが分かる話者だけをテストにした集計を感度として出す．IEMOCAP 即興の 8 つの筋書きは複数のセッションで演じられるので，話題の重なりを限界に書く．
  \item \textbf{発話をそろえる（M7）}：ターゲットごとに，3 モダリティすべての特徴がある発話を 1 つの集合にする．分割，K 本，候補はその上で 1 回だけ作り，7 種の識別器とすべての予測を付ける側で使う．集合の大きさは AFEW 790，IEMOCAP 即興 2,943，MOSEI 18,634，MELD 11,308 である．ソースの識別器も同じ集合で学ぶ．
\end{itemize}

\subsection{値の出所（M8）}

\begin{itemize}
  \item 論文に書く値は，すべて新しい設定の出力から計算する．5.1 節の表の値は使わない．
  \item 中心化に使うターゲットの平均は，テストを除いた候補と K 本から作る．
  \item 説明の解析（7.3 節の順 7）は新しい設定の出力から計算し，読むファイルを設計の文書に挙げる．
\end{itemize}

\subsection{保存の契約（軽微な指摘 m8，m9）}

\begin{itemize}
  \item 全部の段の識別器の重みと切片を，元の特徴の上の値で残す．対象は，コーパスごとの識別器，出発点，予測を付ける側の途中の識別器（K 本で学び直した融合，積の中で学び直した各モダリティの識別器），生徒（2 条件 × 格子 × K × 反復．役割を分けた版は切片を直す前の重みも）である．
  \item 候補の全ラベルで学んだ識別器を，生徒と同じ分割と出発点で作って残す．
  \item 予測を付ける側のすべての種類の確率を，ターゲットの全部の発話について残す．
  \item テスト，K 本，候補，除いた発話の sample\_id と，乱数の種を，反復ごとに残す．
  \item 仮のターゲットの成績を，格子のすべての点について残す．
  \item 凸でない手法（勾配法）は，重みに加えて，テストと候補の確率そのものを残す．
  \item 特徴のファイルの sha256，AFEW の書き起こしの版，コミット，台本の sha256，乱数の種，ライブラリの版，計算機，収束の記録を残す．
\end{itemize}

\subsection{軽微な指摘}

\begin{itemize}
  \item \textbf{m1 メタ学習の語}：この計画の「メタ学習」は，ソースのコーパスを 1 つずつ外した検証で設定を選ぶこと（leave-one-source-out の検証による設定の選択）である．論文ではこの名前で書き，「選択あり・なし」と呼ぶ．規則そのものを学ぶメタ学習は，この計画に入れない．
  \item \textbf{m4 融合の標準化}：融合の特徴は，各モダリティの元の特徴を標準化せずにつなぐ．識別器は自分の学習データで標準化する．全群（テストを含む）の統計による前処理は使わない．読み込みの関数にこの選択肢を足し，検算を流す．
  \item \textbf{m5 計算機}：比較に入る行はすべて RunPod の H100 で回し，行ごとに計算機を記録する．研究室の A5000 で回した部分は，代表の行で H100 と一致するかを確かめる．
  \item \textbf{m6 外すもの}：RAVDESS と CREMA-D は言葉が無く（RAVDESS は顔も無い），AVT を作れないので外す．IEMOCAP の台本は，ターゲットのとき同じ台詞の問題があるので，コーパスの中身をソースとターゲットでそろえるためにソースでも外す．
  \item \textbf{m7 範囲}：結論は凍結した特徴の上の線形の識別器の範囲に限ると書く．AFEW の言葉だけが機械の書き起こしであることを，言葉の識別器の比べに添える．
\end{itemize}

\subsection{研究者に確かめること}

次の既定は，研究者がまだ確かめていない．違えば直す．
\begin{enumerate}
  \item 7.2 節の推す既定（IEMOCAP は即興だけ，資産の定義，C の選び方，平均を引く，重なりの除外，融合の種類）．
  \item 出発点を，コーパスごとの識別器の点数の等しい重みの平均にすること（8.4 節）．
  \item 選択なしの規則を「格子の中央」にすること（8.4 節）．
  \item 主な比較の族 4 つと，同じ水準の幅の規則（8.2 節）．
  \item 反復を「5 つの塊の分割 × K 本の種 2 つ」にし，K を 0，1，5，10，25 にすること（8.2 節）．
\end{enumerate}

\section{用語}

\begin{description}[style=nextline,leftmargin=8mm]
  \item[leave-one-corpus-out] コーパスを 1 つずつターゲットにし，残りをソースにして，全部のコーパスを順に試す評価の方法．
  \item[疑似ラベル] ラベル無しの発話に，ある識別器が付けた予測（感情ごとの確率）．
  \item[候補] ターゲットのラベル無しの発話．疑似ラベルを付ける先である．
  \item[K 本] ターゲットで感情ごとに K 本だけあるラベル付きの発話．4 感情なので 4K 本である．
  \item[出発点] ターゲットで学び始めるときの識別器．離れすぎないように，学び直しのときに L2 正則化で引き戻す先でもある．
  \item[位置の誤り] 感情ごとの点数の底上げ（切片）のずれで失う成績．
  \item[方向の誤り] 重みの向きのずれで失う成績．感情の対ごとに測る．
  \item[仮のターゲット] ソースのコーパスを 1 つ外してターゲットに見立てたもの．設定はここで選び，評価のターゲットのラベルは使わない．
  \item[メタ学習（選択あり）] ソースのコーパスを 1 つずつ外して仮のターゲットにし，その成績で手法の設定を選ぶこと．論文では leave-one-source-out の検証による設定の選択と書く．
  \item[UAR] 感情ごとの正解率を平均した値．でたらめに答えると 25\% になる．
\end{description}

\end{document}
